\section{Proof Sketches}
\label{app:proofs}

\subsection{Proof of Lemma~\ref{lem:first-divergence}}

\begin{proof}
Let $u=(u_1,\ldots,u_m)$ and $v=(v_1,\ldots,v_n)$ be two distinct token
continuations.  If neither sequence is a prefix of the other, then the set
\[
\{t : 1\le t\le \min(m,n),\ u_t\neq v_t\}
\]
is nonempty, and the least element is the first-divergence position.  If one
sequence is a strict prefix of the other, then the first position after the
shorter sequence ends is the first divergence when end-of-sequence is treated
as a token.  Thus every pair of distinct tokenized continuations has a first
divergence after their longest shared token prefix.
\end{proof}

\subsection{Diagnostic Consequence}

The lemma does not state that final text exposes a reliable public verifier.
It only states that any token-level difference has an earliest divergent
event.  This is why first-divergence traces are useful as an upper-bound
diagnostic: if a provider-side event trace cannot recover a signal, then the
weaker public final-text view is unlikely to recover that same signal.  If the
trace recovers a signal while public final text does not, the observed
bottleneck is public observability rather than event-level controllability.
